% Part: second-order-logic
% Chapter: sol-and-sets
% Section: cardinalities

\documentclass[../../../include/open-logic-section]{subfiles}

\begin{document}

\olfileid[pt]{sol}{set}{crd}

\olsection{Cardinalidades de conjuntos}

\begin{explain}
Assim como podemos expressar que o domínio é finito ou infinito,
!!{enumerable} ou !!{nonenumerable}, podemos definir a propriedade de um
subconjunto de~$\Domain{M}$ ser finito ou infinito, !!{enumerable} ou
!!{nonenumerable}.
\end{explain}

\begin{prop}
A fórmula $\fn{Inf}(X) \ident$
\begin{multline*}
\lexists[u][(\lforall[x][\lforall[y][(\eq[u(x)][u(y)] \lif 
      \eq[x][y])]] \land {} \\\lexists[y][(X(y) \land \lforall[x][(X(x)
      \lif \eq/[y][u(x)])]])]
\end{multline*}
é satisfeita em relação a uma atribuição de variáveis~$s$ sse $s(X)$ é
infinito.
\end{prop}

\begin{prop}
A fórmula $\fn{Count}(X) \ident $
\begin{multline*}
\lexists[z][\lexists[u][(X(z) \land 
    \lforall[x][(X(x) \lif X(u(x)))] \land {} \\ \lforall[Y][((Y(z) \land
      \lforall[x][(Y(x) \lif Y(u(x)))]) \lif X = Y])]])
\end{multline*}
é satisfeita em relação a uma atribuição de variáveis~$s$ sse $s(X)$ é
!!{enumerable}.
\end{prop}

Sabemos, pelo teorema de Cantor, que existem conjuntos !!{nonenumerable} e,
de fato, que há infinitos níveis distintos de tamanhos infinitos. A teoria
dos conjuntos desenvolve toda uma aritmética dos tamanhos de conjuntos e
atribui números cardinais infinitos aos conjuntos. Os números naturais
servem como números cardinais que medem os tamanhos dos conjuntos finitos. A
cardinalidade dos conjuntos !!{denumerable} é a primeira cardinalidade
infinita, chamada~$\aleph_0$ (``alef-zero''). O próximo tamanho infinito é
~$\aleph_1$. É o menor tamanho que um conjunto pode ter sem ser contável
(O.s.\@, sem ter tamanho~$\aleph_0$). Podemos definir ``$X$ tem tamanho
$\aleph_0$'' como
$\fn{Aleph}_0(X) \liff \fn{Inf}(X) \land \fn{Count}(X)$. $X$ tem tamanho
$\aleph_1$ sse todos os seus subconjuntos são finitos ou têm tamanho
~$\aleph_0$, mas ele próprio não tem tamanho~$\aleph_0$. Logo, podemos
expressar isso pela fórmula $\fn{Aleph_1}(X) \ident
\lforall[Y][(Y \subseteq X \lif
  (\lnot\fn{Inf}(Y) \lor \fn{Aleph}_0(Y)))] \land \lnot
\fn{Aleph}_0(X)$. Ter tamanho $\aleph_2$ é definido de modo semelhante, e
assim por diante.

Há um tamanho de interesse especial, a chamada cardinalidade do contínuo.
É o tamanho de $\Pow{\Nat}$ ou, equivalentemente, o tamanho de~$\Real$.
Também é possível expressar na lógica de segunda ordem que um conjunto tem
o tamanho do contínuo, mas isso exige um pouco mais de trabalho.

\end{document}
